\documentclass[10pt,a4paper]{amsart}
\usepackage[margin=19mm]{geometry}
\usepackage{amsmath,amssymb,mathtools}
\usepackage[scr=rsfs,cal=euler]{mathalfa}
\usepackage{xfrac}
\usepackage[unicode,hidelinks,bookmarks=false]{hyperref}
\newcommand{\ie}{i.e.\ }
\newcommand{\cf}{cf.\ }
\newcommand{\resp}{respectively\ }
\newcommand{\Subsection}{\subsection}
\providecommand{\nobreakdash}{\nobreak\hbox{-}}
\setlength{\emergencystretch}{2em}
\multlinegap0pt
\usepackage{bm}
\usepackage{bbm}
\usepackage{dsfont}
\usepackage{xspace}
\usepackage{cancel}
\usepackage{mathtools}
\usepackage{cleveref}
\usepackage{amsthm}
\makeatletter
\@ifundefined{theorem}{\newtheorem{theorem}{Theorem}}{}
\@ifundefined{lemma}{\newtheorem{lemma}[theorem]{Lemma}}{}
\@ifundefined{proposition}{\newtheorem{proposition}[theorem]{Proposition}}{}
\@ifundefined{remark}{\newtheorem{remark}[theorem]{Remark}}{}
\makeatother
\def\bV{{\mathbf{V}}}
\def\bm{{\mathsf{m}}}
\def\bq{{\mathsf{q}}}
\def\i{{(\iota)}}
\def\sU{{\mathsf{U}}}
\def\sV{{\mathsf{V}}}
\newcommand{\Dx}{{\Delta x}}
\newcommand{\D}{\Delta}
\newcommand{\NN}{{\mathbb{N}}}
\newcommand{\N}{\mathbb{N}}
\newcommand{\R}{\mathbb{R}}
\newcommand{\bc}{\mathbf{c}}
\newcommand{\beq}{\begin{equation}}
\newcommand{\eeq}{\end{equation}}
\newcommand{\mA}{\mathcal A}
\newcommand{\mC}{\mathcal C}
\newcommand{\ux}{{\underline x}}
\title[A semi-Lagrangian method for solving state constraint Mean F]{A semi-Lagrangian method for solving state constraint Mean Field Games in Macroeconomics}
\author{Fabio Camilli, Qing Tang, Yong-shen Zhou}
\date{Source: arXiv:2510.00768v1 (CC BY 4.0)}
\begin{document}
\maketitle

\noindent\emph{Source and licence.} A semi-Lagrangian method for solving state constraint Mean Field Games in Macroeconomics. Version arXiv:2510.00768v1 (CC BY 4.0), distributed under the Creative Commons Attribution 4.0 International licence, as recorded in the arXiv licence metadata for this exact version and shown on the abstract page https://arxiv.org/abs/2510.00768v1. Full rights notice: licenses/arxiv-2510.00768-CC-BY-4.0.txt.\par\smallskip

\noindent\textit{Editorial scope.} Counted unit: the Proposition whose source label is \texttt{prop\_scheme} in the article (source0.tex lines 1171--1185, source label \texttt{prop\_scheme}), delivered together with the article's own complete original proof of it (source0.tex lines 1186--1265, one \texttt{proof} environment of 80 lines). No mathematics is written, repaired or re-derived: statement and proof are the article's own text line for line. Required context retained uncounted: C set (lines 766--768); Scheme (lines 971--973); def S (lines 975--981); C\_j (lines 983--989); monotone\_scheme (lines 1032--1038); S monotone (lines 1035--1037); M pos (lines 1039--1067); U leq V (lines 1039--1067); eq:policy\_ev (lines 1115--1117); update\_c (lines 1120--1126); Howard theo (lines 1128--1134). Every definition, display and label of those lines is retained unchanged. Omitted from this package and not redistributed: 1--765, 769--970, 974--974, 982--982, 990--1031, 1038--1038, 1068--1114, 1118--1119, 1127--1127, 1135--1170, 1266--1548. Nothing inside a retained excerpt is omitted or altered: every retained body is delivered exactly as the article's lines are written, so no omission receipt of any kind accompanies this package. Citations: the retained excerpts carry no citation command at all, so no bibliography entry of the article is transcribed and no reference is redistributed. Reference adaptation: none was needed and none was made; every reference inside the delivered unit resolves inside it, so no unresolved reference or citation is produced. Printed slips kept literal: no printed word, symbol, hypothesis or number of the counted statement or of its proof was altered. Layout and class: the article's own class is \texttt{article} with packages including amssymb, amsbsy, amsmath, amsfonts, amscd, amsthm, fullpage, fontenc, appendix, hyperref, url, courier, easyReview, cleveref; the wrapper is the lane's pinned \texttt{amsart} editorial wrapper at 10pt on A4, which restates the theorem-like environments the retained text needs and adds the article's own macro definitions that the retained text needs (\texttt{\textbackslash D}, \texttt{\textbackslash Dx}, \texttt{\textbackslash N}, \texttt{\textbackslash NN}, \texttt{\textbackslash R}, \texttt{\textbackslash bV}, \texttt{\textbackslash bc}, \texttt{\textbackslash beq}, \texttt{\textbackslash bm}, \texttt{\textbackslash bq}, \texttt{\textbackslash eeq}, \texttt{\textbackslash i}, \texttt{\textbackslash mA}, \texttt{\textbackslash mC}, \texttt{\textbackslash sU}, \texttt{\textbackslash sV}, \texttt{\textbackslash ux}), with extra wrapper packages (bm, bbm, dsfont, xspace, cancel, mathtools, cleveref). The delivered numbering is the unit's own. Assets: no retained span contains a figure environment, a graphics-inclusion command, a PDF-page inclusion, a table or a TikZ picture; no image file is copied into this package, the package declares no asset dependency and no image file is redistributed with it. Licence: version arXiv:2510.00768v1 of this article is distributed under the Creative Commons Attribution 4.0 International licence, as recorded in the arXiv licence metadata for this exact version and shown on the abstract page https://arxiv.org/abs/2510.00768v1. A full rights notice is delivered at licenses/arxiv-2510.00768-CC-BY-4.0.txt. Characters: every retained excerpt is pure ASCII, so the delivered engine, XeLaTeX with its default Latin Modern text font, reports no missing character.
\begin{equation}\label{C set}
\mC^{h}_{j}(x):=\left\{{c}: c\ge 0\quad\text{and}\quad x+h(rx+y_j-c) \ge \ux \right\}.
\end{equation}

\begin{equation}\label{Scheme}
	\mathcal{F}^\D_j\left(x_i,\left[V^{\Delta}_{i,j},V^{\Delta}_{i,\bar \jmath}\right],V^{\Delta}_{\cdot,j}\right)=0\qquad i\in\NN,\,j=1,2,\,\bar{\jmath}=3-j.
\end{equation}

\begin{equation}\label{def S}
\begin{aligned}
\mathcal{F}^\D_j(x_i,(\bq_j,\bq_{\bar \jmath}),\sU)={}&\rho \bq_j-(1-\rho h)\lambda_j(\bq_{\bar \jmath}-\bq_j)\\
-\sup_{c\in \mC^{\D}_{j}(x_i)}&\left\{
 u(c)+(1-\rho h)(1-\lambda_jh)\frac{1}{h}\left(\sum_k \beta_k(x_i+hs_{i,1}(c))\sU_k-\bq_j  \right)\right\},
\end{aligned}
\end{equation}

\begin{remark}\label{C_j}
We will use the condition for step size $\Dx\sim h$ for $\D\to 0$. Notice in the fully discretized setting \cref{C set} becomes  
\beq\label{Eq:C_j}
\mC^{\D}_{j}(x_i)=\left[0, \frac{x_i-\ux}{h}+rx_i+y_j\right].
\eeq
By imposing $\Dx\sim h$, the constraint is binding only at $\ux$. For any $x>\ux$, $x=i(\D x)\Dx$ where $i(\D x)\to +\infty$ as $\D x\to 0$. Therefore, for any given $R>0$, if $\D$ is sufficiently small and $x_i>\ux$ then $[0,R)\subset \mC^{\D}_{j}(x_i)$. 
\end{remark}

\begin{lemma}\label{monotone_scheme}
For any  $\D$, $i\in\N$,
	bounded functions $\sU$, $\sV\in B(\mA^\Dx)$ such that $\sU_k\le \sV_k$  $\forall k\in\N$ and  $(\bq_1,\bq_2)$, $(\bm_1,\bm_2)\in\R^2$ such that $\displaystyle \theta:=\bq_j-\bm_j=\max_{k=1,2}\{\bq_k-\bm_k\}\ge 0$, then
\begin{equation}\label{S monotone}
\mathcal{F}^\D_j(x_i,(\bq_j,\bq_{\bar \jmath}), \mathsf{U}+\theta)-\mathcal{F}^\D_j(i,(\bm_j,\bm_{\bar \jmath}),\mathsf{V})\ge \rho\theta.
\end{equation}
\end{lemma}

\begin{proof}
%%%%% monotonicity %%%%%
We recall that for all $i$, $j=1,2$ and $c \in \mC^{\D}_{j}$,
\beq\label{M pos}
\beta_k(x_i+hs_{i,1}(c))\ge 0, \quad \sum_k \beta_k(x_i+hs_{i,1}(c))=1.
\eeq
Assume that $\theta:=\bq_1-\bm_1\ge \bq_2-\bm_2$, hence $\bq_2-\bq_1\le \bm_2-\bm_1$. We first obtain from $\sU_k\le \sV_k$ for all $k$ and \cref{M pos} that
\beq\label{U leq V}
\begin{aligned}
&\sup_{c\in \mC^{\D}_{j}(x_i)}\left\{
	u(c)+(1-\rho h)(1-\lambda_1h)\frac{1}{h}\left(\sum_k \beta_k(x_i+hs_{i,1}(c))\sU_k-\bm_1\right)\right\}\\
\leq {}&
\sup_{c\in \mC^{\D}_{j}(x_i)}\left\{
	u(c)+(1-\rho h)(1-\lambda_1h)\frac{1}{h}\left(\sum_k \beta_k(x_i+hs_{i,1}(c))\sV_k-\bm_1\right)\right\}.
\end{aligned}
\eeq
By \cref{def S}, $\bq_1=\bm_1+\theta$ and \cref{M pos} we have
\begin{align*}
&\quad \mathcal{F}^\D_1(x_i,(\bq_1,\bq_2), \sU+\theta)\\
&=\rho(\bm_1+\theta)-(1-\rho h)\lambda_1(\bq_2-\bq_1)\\
       &-\sup_{\bc\in \mC^{\D}_{j}(x_i)}\left\{
	u(c)+\frac{(1-\rho h)(1-\lambda_1h)}{h}\left(\sum_k \beta_k(x_i+hs_{i,1}(c))(\sU_k+\theta)-(\bm_1+\theta)  \right)\right\}\\
	&\geq \rho \theta+\rho \bm_1-(1-\rho h)\lambda_1(\bm_2-\bm_1)\\
       &-\sup_{c\in \mC^{\D}_{j}(x_i)}\left\{
	u(c)+\frac{(1-\rho h)(1-\lambda_1h)}{h}\left(\sum_k \beta_k(x_i+hs_{i,1}(c))\sU_k-\bm_1\right)\right\}.
\end{align*}
We can then apply \cref{U leq V} to obtain \cref{S monotone}. We proceed similarly if $\theta:=\bq_2-\bm_2\ge \bq_1-\bm_1$. Hence monotonicity of the scheme is proved.\par
%%%%% stability %%%%%% 
 \end{proof}

\begin{equation}\label{eq:policy_ev}
     V_{i,j}^{\Delta, \i} = hu(c_{i,j}^{\i})+(1-\rho h)\left[\lambda_jhV_{i,\bar{\jmath}}^{\Delta,\i}+(1-\lambda_jh)\left(\sum_k\beta_k\left(x_i+hs^\i_{i,j}\right)V_{k,j}^{\Delta,\i}\right)\right].
\end{equation}\\

\begin{equation}\label{update_c}
\begin{aligned}
        &c^{(\iota+1)}_{i,j}=\mathop{\arg\max}\limits_{c\in \mC^{\D}_{j}(x_i)}\left\{
		hu(c)+(1-\rho h)(1-\lambda_jh)\left(\sum_k \beta_k(x_i+hs_{i,j}(c))V^{\D,\i}_{k,j}\right)\right\},\\
&s^{(\iota+1)}_{i,j}=s_{i,j}(c^{(\iota+1)}).		
\end{aligned}
\end{equation}

\begin{theorem}\label{Howard theo}
	Let $\bV^\i=(\bV^\i_1, \bV^\i_2)$, $\bV^\i_j\in  B(\mA^\Dx)$, be the sequence generated by the policy iteration method \eqref{eq:policy_ev}-\eqref{update_c}. Then
	\begin{equation}\label{eq:est_PI}
		\bV^\i \le \bV^{(\iota+1)} \qquad \forall \iota \in \N.
	\end{equation}
Moreover, $\lim_{\iota \to \infty}\bV^\i=\bV$ such that $\bV=(\bV_1,\bV_2)\in B(\mA^\Dx)^2$ is the unique solution of \cref{Scheme}.
\end{theorem}

\begin{proposition}\label{prop_scheme}
Assume that $\D x\sim h$ for $\D\to 0$. The scheme $\mathcal{F}^\D$, besides the monotonicity property established in \cref{monotone_scheme}, satisfies the following properties:
\begin{itemize}
\item[]{\em  Stability: } The unique  solution $\bV=(V^\D_{\cdot,1},V^\D_{\cdot,2})$ to \cref{Scheme} is uniformly bounded in $\D$.
	
	\item[]{\em  Consistency: } For $j=1,2$ and for any
	smooth function $\phi=(\phi_1,\phi_2)$
	\begin{align*}
			\limsup_{{\D\to 0}\atop{x_i\to x} } \mathcal{F}^\D_j(x_i,\phi(x_i),\phi_{j})  & \leq  \rho \phi_{j}(x)- H(x,y_j,D\phi_{j})-\lambda_j(\phi_{j}(x)-\phi_{\bar \jmath}(x))
		 \quad \forall x\in(\ux,\infty),\\
		\liminf_{{\D\to 0}\atop{x_i\to x}} \mathcal{F}^\D_j(x_i,\phi(x_i),\phi_{j})    &\geq \rho \phi_{j}(x)- H(x,y_j,D\phi_{j})-\lambda_j(\phi_{j}(x)-\phi_{\bar \jmath}(x)) 
	    \quad \forall x\in[\ux,\infty).
	\end{align*}
	\end{itemize}
\end{proposition}

\begin{proof}
The existence and uniqueness of solution to \cref{Scheme} has been considered in \cref{Howard theo}. We now use comparison to give the uniform bound. Clearly, $(0,0)$ is a supersolution to \cref{Scheme}. We now show the constant
\beq\label{def discrete sub}
\left(\frac{1}{\rho}\frac{(r\ux+y_1)^{1-\gamma}}{1-\gamma},\frac{1}{\rho}\frac{(r\ux+y_1)^{1-\gamma}}{1-\gamma}\right)
\eeq
 is a subsolution. Since \cref{def discrete sub} is a constant solution and from \cref{M pos}, we only need to show
\beq\label{discrete sub}
\frac{(r\ux+y_1)^{1-\gamma}}{1-\gamma}\leq \sup_{c\in \mC^{\D}_{j}(x_i)}\{u(c)\}\quad \forall x_i.
\eeq
From \cref{C set}, $\sup_{c\in \mC^{\D}_{j}(\ux)}\{u(c)\}=\frac{(r\ux+y_j)^{1-\gamma}}{1-\gamma}$ and $\sup_{c\in \mC^{\D}_{j}(x_i)}\{u(c)\}\geq \sup_{c\in \mC^{\D}_{j}(\ux)}\{u(c)\}$ if $x_i>\ux$ . We conclude \cref{discrete sub} by observing $y_1<y_2$.  \par
We now consider the consistency of scheme.  Given a  function $\phi=(\phi^1,\phi^2)$ with $\phi_{j}:\R\to\R$, we denote with 
$
	I[\phi_{j}](x):=\sum_{i\in  \N} \beta_i(x) \phi_{j} (x_i) 
$
its linear interpolation on the grid $\mA^\Dx$. If $\phi_{j}\in C^2(\R)$ with bounded derivative, then there exists a positive constant $C_1$ such that
\begin{equation}\label{eq:intepolation} 
	\sup_{x\in \R}| I[ \phi_{j}](x)-\phi_{j}(x)|\le C_1(\Dx)^2.
\end{equation} 
From \eqref{eq:intepolation} we obtain
\begin{equation}\label{eq:estimate_grad}
	\begin{split}
		& \left|\frac{1}{h} \big[ \sum_k \beta_k(x_i+hs_{i,j}(c))\phi_{j}(x_i)-\phi_{j}(x_i)]-D\phi_{j}(x_i)s_j(x_i)\right|\\
		=& \left|\frac{1}{h} \big[I[ \phi_{j}](x_i+hs_{i,j}(c))-\phi_{j}(x_i)\big]-D\phi_{j}(x_i)s_j(x_i)\right| \\
\le &\left|\frac{1}{h}  \big[\phi_{j}(x_i+hs_{i,j}(c))-\phi_{j}(x_i)\big]-D\phi_{j}(x_i)s_j(x_i)\right| +	C_1\frac{(\Dx)^2}{h}
\le C\left(\frac{(\Dx)^2}{h}+h\right),
	\end{split}
\end{equation}
we have from \cref{eq:estimate_grad}:
\begin{align*}
&\sup_{c\in \mC^{\D}_{j}(x_i)}\left\{
	u(c)+(1-\rho h)(1-\lambda_jh)D\phi_{j}(x_i)(rx_i+y_j-c)
	\right\}\\
\leq {}&\sup_{c\in \mC^{\D}_{j}(x_i)}\left\{
	u(c)+(1-\rho h)(1-\lambda_jh)\frac{1}{h} \left[ \sum_k \beta_k\left(x_i+hs_{i,j}(c)\right)\phi_{j}(x_i)-\phi_{j}(x_i)\right] 
	\right\}
	+C\left(\frac{(\Dx)^2}{h}+h\right).
\end{align*}
Let $x_i\to x\in (\ux,\infty)$ for $\D \to 0$, then $x_i>\ux$ when $\D$ is sufficiently small. From \cref{C_j}, we have
$
H(x_i, y_j,D\phi_{j})=\sup_{c\in \mC^{h}_{j}(x_i)}\left\{
	u(c)+(rx_i+y_j-c)D\phi_{j}(x_i)\right\} .
$
We then obtain
\begin{align*}
	  &\rho \phi_{j}(x_i)- H(x_i,y_j,D\phi_{j})-\lambda_j(\phi^{j}(x_i)-\phi_{\bar \jmath}(x_i))\\
	={}&\rho \phi_{j}(x_i)-\sup_{c\in \mC^{h}_{j}(x_i)}\left\{
	u(c)+(rx_i+y_j-c)D\phi_{j}(x_i)\right\} - \lambda_j(\phi_{\bar \jmath}(x_i)-\phi_{j}(x_i)) \\
	\ge{}& \rho \phi_{j}(x_i)-\sup_{c\in \mC^{h}_{j}(x_i)}\left\{
	u(c)+(1-\rho h)(1-\lambda_jh)D\phi_{j}(x_i)(rx_i+y_j-c)
	\right\}\\
	&-(1-\rho h)\lambda_j(\phi_{\bar \jmath}(x_i)-\phi_{j}(x_i))-C_1h\\
	\ge{}& \rho \phi_{j}(x_i)-(1-\rho h)\lambda_j(\phi_{\bar \jmath}(x_i)-\phi_{j}(x_i))-C\left(\frac{(\Dx)^2}{h}+h\right)\\
	&-\sup_{c\in \mC^{h}_{j}(x_i)}\left\{
	u(c)+\frac{(1-\rho h)(1-\lambda_jh)}{h} \left[ \sum_k \beta_k(x_i+hs_{i,j}(c))\phi_{j}(x_i)-\phi_{j}(x_i)\right] 
	\right\}
\end{align*}
and passing to the $\limsup$ in $(\ux,\infty)$,  we get the first condition in
{\em consistency}. 

%
% for   $L=\sup\|D\phi\|_\infty$, we have
%	\begin{equation}\label{eq:Lipschitz} 
%		\sup_{a\in \R}| I[ \phi](x)-\phi(x)|\le L \Dx.
%	\end{equation}  
Given $x\in [\ux,\infty)$, let $x_i\in \mA^\Dx$ such that $x_i\to x$. Then, taking into account \eqref{eq:estimate_grad}, we have
\begin{align*}
&\mathcal{F}^\D_j(x_i,\phi(x_i),\phi_{j})\\
\ge{}& \rho \phi_{j}(x_i)-\sup_{c\in \mC^{\D}_{j}(x_i)}\left\{
	u(c)+(1-\rho h)(1-\lambda_jh)D\phi_{j}(x_i)(rx_i+y_j-c)
	\right\}\\
	&-(1-\rho h)\lambda_j(\phi_{\bar \jmath}(x_i)-\phi_{j}(x_i))-C_1\left(h+\frac{(\Dx)^2}{h}\right)\\
\ge{}& \rho \phi_{j}(x_i)-\sup_{c\in (0,\infty)}\Big\{
	u(c)+(1-\rho h)(1-\lambda_jh)D\phi_{j}(x_i)(rx_i+y_j-c)
	\Big\}\\
	&-(1-\rho h)\lambda_j(\phi_{\bar \jmath}(x_i)-\phi_{j}(x_i))-C_1\left(h+\frac{(\Dx)^2}{h}\right)\\
={}& \rho \phi_{j}(x_i)- H(x_i, y_j,D\phi_{j})-\lambda_j(\phi_{j}(x_i)-\phi_{\bar \jmath}(x_i))-C\left(\frac{(\Dx)^2}{h}+h\right).
\end{align*}
Passing to the $\liminf_{{\D\to 0}\atop{x_i\to x}}$ in the previous inequality, we get the second condition in {\em consistency}.

\end{proof}

\end{document}
